% !TEX root = ../main.tex
% !TEX program = lualatex
\documentclass[../main.tex]{subfiles}
\IfSubfilesClassLoaded{\externaldocument{\subfix{../build/main}}}{}
% =============================================================================
% CHAPTER 10: COST ESTIMATING
% DESCRIPTION: Life-cycle cost scope, methods, products/players, and learning curves.
% =============================================================================
\begin{document}
\IfSubfilesClassLoaded{\chapter{Cost Estimating}}{}
%====================
% Contents here
%====================

% === COST ESTIMATING ===
\section{Cost Estimating}

\subsection{Summary}
\begin{description}
  \item[\textbf{Cost scope.}] \ac{lcc} includes acquisition, operating and support, and disposal costs; \ac{toc}
  broadens the view by adding enterprise support, infrastructure, and indirect ownership costs.
  \item[\textbf{Estimate methods.}] Use analogy and parametric methods early, engineering/bottom-up estimates when the
  design is detailed, and actuals when production or sustainment data exist.
  \item[\textbf{Decision products.}] The \ac{pmo} owns the program definition and \ac{card}; the \ac{poe}, \ac{cce},
  \ac{ccp}, \ac{ice}, and CAPE review lanes test whether the cost position is realistic enough for milestone and
  budget decisions.
\end{description}
~\srcCite{edo-3-1-6-cost-estimating-2024}

\subsection{Life-cycle cost composition}
\begin{description}
  \item[\textbf{Flyaway/Sailaway/Rollaway cost.}] Cost of procuring only the prime mission equipment (aircraft, ship, tank)~\autocite{edo-3-1-6-cost-estimating-2024}.
  \item[\textbf{Weapon system cost.}] Prime mission equipment plus support items required to deliver a functioning system (e.g., manuals, test equipment)~\autocite{edo-3-1-6-cost-estimating-2024}.
  \item[\textbf{Procurement cost.}] Weapon system cost plus support items and initial spares required to stock the supply system~\autocite{edo-3-1-6-cost-estimating-2024}.
  \item[\textbf{Program acquisition cost.}] Costs to develop, procure, and house the system, including facilities~\autocite{edo-3-1-6-cost-estimating-2024}.
  \item[\textbf{Operating and support cost.}] All costs to operate, modify, maintain, and support the system in the inventory~\autocite{edo-3-1-6-cost-estimating-2024}.
  \item[\textbf{Disposal cost.}] Demilitarization and disposal after useful life~\autocite{edo-3-1-6-cost-estimating-2024}.
\end{description}

\note{\ac{lcc} equals program acquisition cost plus operating and support plus disposal; \ac{os.phase} is typically the largest share (about two-thirds) for most programs~\autocite{edo-3-1-6-cost-estimating-2024}.}

\begin{figure}[H]
  \centering
  \includegraphics[width=0.7\linewidth]{LCC_Composition.png}
  \caption[Life-cycle cost composition]{Life-cycle cost composition. \srcCite{edo-3-1-6-cost-estimating-2024}}
  \label{fig:lcc_composition}
\end{figure}

\begin{figure}[H]
  \centering
  \includegraphics[width=0.7\linewidth]{lcc_over_time.png}
  \caption[Life-cycle cost over time]{Notional life-cycle cost distribution over time. \srcCite{edo-3-1-6-cost-estimating-2024}}
  \label{fig:lcc_time}
\end{figure}

\subsection{Total ownership cost (TOC)}
\begin{description}
  \item[\textbf{TOC.}] \ac{toc} includes \ac{lcc} plus common spares/support systems, infrastructure, and distributed indirect costs across the \ac{dod} enterprise~\autocite{edo-3-1-6-cost-estimating-2024}.
  \item[\textbf{TOC vs LCC.}] \ac{lcc} is a subset of \ac{toc} and focuses on direct program costs; \ac{toc} seeks to allocate additional indirect and infrastructure costs to the program~\autocite{edo-3-1-6-cost-estimating-2024}.
\end{description}

\subsection{Cost estimating methods}
\begin{description}
  \item[\textbf{Analogy.}] Compares to a similar system and adjusts for differences; used very early when data are limited~\autocite{edo-3-1-6-cost-estimating-2024}.
  \item[\textbf{Parametric (CER).}] Uses statistical cost-estimating relationships (e.g., weight, speed, schedule) when key parameters are known~\autocite{edo-3-1-6-cost-estimating-2024}.
  \item[\textbf{Engineering (bottom-up).}] Builds from a detailed \ac{wbs}; most accurate but time-consuming and data-intensive~\autocite{edo-3-1-6-cost-estimating-2024}.
  \item[\textbf{Actuals (extrapolation).}] Uses actual costs from the same system or close predecessor; strongest in production~\autocite{edo-3-1-6-cost-estimating-2024}.
\end{description}

\begin{longtblr}[
    label   = {tab:cost_method_phases},
    caption = {Cost estimating methods by acquisition phase},
    entry   = {Estimating Methods by Phase},
    remark{\tSource} = {\tabCite{edo-3-1-6-cost-estimating-2024}}
  ]
  {colspec = {@{} ll @{}}}
  \toprule
  {Phase} & {Most applicable methods} \\
  \midrule
  MSA  & Analogy, Parametric \\
  TMRR & Parametric, Engineering \\
  EMD  & Engineering, Actuals \\
  P\&D & Engineering, Actuals \\
  O\&S & Actuals \\
  \bottomrule
\end{longtblr}

\note{Pool example: analogy (neighbor costs), parametric (cost per square foot), engineering (detailed materials and labor), actuals (later builds in a housing development)~\autocite{edo-3-1-6-cost-estimating-2024}.}

\subsection{Cost estimating products and players}
\begin{longtblr}[
    label   = {tab:cost_players_products},
    caption = {Cost estimating players and products},
    entry   = {Cost Players and Products},
    remark{\tSource} = {\tabCite{edo-3-1-6-cost-estimating-2024}}
  ]
  {colspec = {@{} ll @{}}}
  \toprule
  {Player} & {Product} \\
  \midrule
  Program Management Office / SYSCOM cost organization & CARD, economic analysis, POE \\
  Component cost agency (e.g., NCCA) & ICE, CCE, CCP \\
  CAPE & ICE \\
  \bottomrule
\end{longtblr}

\begin{description}
  \item[\textbf{Program management office.}] Defines the program, owns the CARD and program estimate need, and may use the SYSCOM cost organization to develop the POE on behalf of the PMO~\autocite{edo-3-1-6-cost-estimating-2024}.
  \item[\textbf{Component cost agency.}] Leads cost analysis independent of the acquisition decision chain; prepares \ac{cce} for ACAT IC/ID programs and may serve as the \ac{ice}~\autocite{edo-3-1-6-cost-estimating-2024}.
  \item[\textbf{CAPE.}] Conducts or approves \acp{ice} for ACAT ID/IC programs, participates in PPBE reviews, and advises \ac{secdef}~\autocite{edo-3-1-6-cost-estimating-2024}.
\end{description}

\subsection{Cost estimating products}
\begin{description}
  \item[\textbf{CARD.}] Developed by the program office; provides quantitative technical and programmatic descriptions for costing; updated for all milestone reviews (\ac{dod} Manual 5000.04 series)~\autocite{edo-3-1-6-cost-estimating-2024}.
  \item[\textbf{CCE.}] Developed by the component cost agency using input from the POE; reviewed by CAPE for \acp{mdap} and \acp{mais}~\autocite{edo-3-1-6-cost-estimating-2024}.
  \item[\textbf{POE.}] Program estimate developed by or for the \ac{pmo}, often by the SYSCOM cost organization; estimates full life-cycle resources and costs based on the CARD~\autocite{edo-3-1-6-cost-estimating-2024}.
  \item[\textbf{CCP.}] Reconciliation between the CCE and POE; official component cost position (DoN typically uses the POE in lieu of a CCP)~\autocite{edo-3-1-6-cost-estimating-2024}.
  \item[\textbf{ICE.}] Independent \ac{lcc} required for ACAT I programs; conducted by an organization independent of the \ac{pmo}~\autocite{edo-3-1-6-cost-estimating-2024}.
  \item[\textbf{Economic analysis.}] Evaluates relative costs of technical alternatives and acquisition strategies; developed by the \ac{pmo} for \acp{mais}~\autocite{edo-3-1-6-cost-estimating-2024}.
\end{description}

\subsection{Typical CARD elements}
\begin{longtblr}[
    label   = {tab:card_elements},
    caption = {Typical CARD elements},
    entry   = {CARD Elements},
    remark{\tSource} = {\tabCite{edo-3-1-6-cost-estimating-2024}}
  ]
  {colspec = {@{} ll @{}}}
  \toprule
  {System and program description} & {Support and planning} \\
  \midrule
  System description and WBS & System logistics and support concepts \\
  Subsystem descriptions and tech maturity & Hardware/software maintenance and support \\
  Quality factors and R/M/A & Training concept and manpower requirements \\
  Risk assessment and mitigation & Time-phased quantities and activity rates (OPTEMPO) \\
  Operational concept and basing & Milestone schedule and acquisition strategy \\
  \bottomrule
\end{longtblr}

\subsection{Learning curve theory}
\begin{description}
  \item[\textbf{Definition.}] When production quantity doubles, per-unit labor hours decrease by a constant percentage~\autocite{edo-3-1-6-cost-estimating-2024}.
  \item[\textbf{Drivers.}] Experience, uninterrupted production, consistent design, tooling improvements, and productivity focus~\autocite{edo-3-1-6-cost-estimating-2024}.
  \item[\textbf{Rule.}] Learning curves apply to production labor only and to repeatable processes (e.g., a 90\% curve makes unit 4 labor 90\% of unit 2)~\autocite{edo-3-1-6-cost-estimating-2024}.
\end{description}

\subsection{Escalation}
\begin{description}
  \item[\textbf{Purpose.}] Escalation accounts for inflation and outlay rates to convert constant-year estimates into then-year budget dollars~\autocite{edo-3-1-6-cost-estimating-2024}.
  \item[\textbf{Indices.}] Raw (compound) index uses inflation only; weighted (composite) index combines inflation and outlay rates and is unique to each appropriation~\autocite{edo-3-1-6-cost-estimating-2024}.
  \item[\textbf{Source.}] \ac{osd} provides annual inflation and outlay rates; services compute indices for budget submissions~\autocite{edo-3-1-6-cost-estimating-2024}.
\end{description}

\subsection{Will cost / should cost}
\begin{description}
  \item[\textbf{Will cost.}] Funded-to level without additional cost-saving initiatives~\autocite{edo-3-1-6-cost-estimating-2024}.
  \item[\textbf{Should cost.}] Desired cost after management initiatives; each savings opportunity must be tied to a quantifiable engineering or business change~\autocite{edo-3-1-6-cost-estimating-2024}.
  \item[\textbf{Opportunities.}] Learning curve gains, process improvements, contracting changes, competition, incentives, and \ac{gfe}~\autocite{edo-3-1-6-cost-estimating-2024}.
  \item[\textbf{Avoid.}] Straight percentage cuts or large investments with no near-term recovery~\autocite{edo-3-1-6-cost-estimating-2024}.
\end{description}

\subsection{Coursebook cue: cost-estimate review}

Coursebook 3.1.6 emphasizes that a cost estimate is not complete just because
the math is finished.  The PMO prepares and defends the program estimate, the
SYSCOM or component cost organization reviews the estimate for realism and
method, and the CAPE role becomes central for major programs that require
independent cost assessment before a milestone decision
\autocite{edo-3-1-6-cost-estimating-2024}.

\begin{longtblr}[
  caption = {Cost-estimate products and review roles.},
  label = {tab:cost-estimate-review-roles},
]{
  width = \linewidth,
  colspec = {X[0.65,l] X[1.0,l] X[1.35,l]},
  rowhead = 1,
  hlines,
  vlines,
}
\textbf{Product} & \textbf{Typical owner} & \textbf{Board cue} \\
CARD & Program office & Describes the program, technical baseline, schedule, quantities, and assumptions that estimators price. \\
POE & PMO / SYSCOM cost organization & Program estimate used to plan, budget, and defend the program. \\
CCE & Component cost agency & Independent component-level estimate or assessment used to challenge the program view. \\
ICE & Independent cost authority & Required when statute, policy, acquisition category, or decision authority requires an independent estimate. \\
CAPE review & OSD CAPE & Provides independent cost analysis for major decisions and helps the MDA judge affordability and realism. \\
\end{longtblr}

\subsection{Quick Review}
\begin{itemize}
  \item What is \ac{lcc}? \textbf{Answer:} Acquisition cost plus operating and support plus disposal.
  \item Which estimating method fits which phase? \textbf{Answer:} Early concept uses analogy; preliminary design uses
  parametric; detailed design uses engineering; production and sustainment use actuals.
  \item What is the learning-curve trap? \textbf{Answer:} Learning curves apply to repeatable production labor, not
  material.
  \item What does escalation do? \textbf{Answer:} Converts constant-year estimates into then-year budget dollars using
  inflation and outlay rates.
  \item What makes a defensible estimate? \textbf{Answer:} A technical baseline, documented assumptions, independent
  review, and a clear decision use; know who owns the \ac{card}, \ac{poe}, component review, \ac{ice}, and CAPE role.
\end{itemize}

%====================
% End of file
%====================
\IfSubfilesClassLoaded{
  \RenewDocumentCommand{\entryname}{}{\textbf{\color{Modern} Acronym}}
  \RenewDocumentCommand{\descriptionname}{}{\textbf{\color{Modern} Definition}}
  \printnoidxglossary[
    type=\acronymtype,
    title=Acronyms,
    style=long-booktabs
  ]}{}

\end{document}
